% Created 2024-01-08 lun 09:13
% Intended LaTeX compiler: pdflatex
\documentclass[11pt]{article}
\usepackage[utf8]{inputenc}
\usepackage[T1]{fontenc}
\usepackage{graphicx}
\usepackage{longtable}
\usepackage{wrapfig}
\usepackage{rotating}
\usepackage[normalem]{ulem}
\usepackage{amsmath}
\usepackage{amssymb}
\usepackage{capt-of}
\usepackage{hyperref}
\usepackage{minted}

\usepackage{parskip}
\usepackage[utf8]{inputenc} %% For unicode chars
\usepackage{placeins}
\usepackage[margin=2.50cm]{geometry}
\usepackage[T1]{fontenc}
\usepackage{mathpazo}
\linespread{1.05}
\usepackage[scaled]{helvet}
\usepackage{courier}
\hypersetup{colorlinks=true,linkcolor=blue}
\RequirePackage{fancyvrb}
\DefineVerbatimEnvironment{verbatim}{Verbatim}{fontsize=\small,formatcom = {\color[rgb]{0.5,0,0}}}
\usepackage{mdframed}
\BeforeBeginEnvironment{minted}{\begin{mdframed}}
\AfterEndEnvironment{minted}{\end{mdframed}}
\setcounter{secnumdepth}{3}
\author{José Manuel Sánchez Álvarez}
\date{}
\title{PHP CRUD Car Dealership}
\hypersetup{
 pdfauthor={José Manuel Sánchez Álvarez},
 pdftitle={PHP CRUD Car Dealership},
 pdfkeywords={},
 pdfsubject={},
 pdfcreator={Emacs 28.1 (Org mode 9.5.5)}, 
 pdflang={Spanish}}
\begin{document}

\maketitle
\setcounter{tocdepth}{3}
\tableofcontents

\setlength\parindent{10pt}

\section{Car Dealership CRUD in PHP}
\label{sec:org5a9b2ab}
Creating a CRUD (Create, Read, Update, Delete) application for a car
dealership in PHP involves designing a suitable database schema. Below,
I'll outline a basic database structure with tables, attributes, types,
constraints, keys, and relationships. This schema is designed to manage
information about cars, brands, sellers, and clients.

\subsection{Database}
\label{sec:org281fe6d}
\subsubsection{Table: \texttt{CarBrand}}
\label{sec:org53ed497}
\begin{itemize}
\item \textbf{Attributes:}

\begin{itemize}
\item \texttt{id} (INT, primary key, auto-increment)
\item \texttt{name} (VARCHAR, unique, not null)
\end{itemize}

\item \textbf{Primary Key:} \texttt{id}
\item \textbf{Unique Key:} \texttt{name}
\item \textbf{Purpose:} Stores different car brands.
\end{itemize}

\subsubsection{Table: \texttt{Seller}}
\label{sec:orgd161478}
\begin{itemize}
\item \textbf{Attributes:}

\begin{itemize}
\item \texttt{id} (INT, primary key, auto-increment)
\item \texttt{name} (VARCHAR, not null)
\item \texttt{cars\_selled} (INT, default 0)
\end{itemize}

\item \textbf{Primary Key:} \texttt{id}
\item \textbf{Purpose:} Stores information about car sellers.
\end{itemize}

\subsubsection{Table: \texttt{Car}}
\label{sec:orge1f9115}
\begin{itemize}
\item \textbf{Attributes:}

\begin{itemize}
\item \texttt{id} (INT, primary key, auto-increment)
\item \texttt{model} (VARCHAR, not null)
\item \texttt{price} (DECIMAL(10,2), not null)
\item \texttt{condition} (VARCHAR, not null, values like 'new', 'refurbished')
\item \texttt{engine\_features} (TEXT, nullable)
\item \texttt{gas\_type} (VARCHAR, nullable)
\item \texttt{mileage} (INT, nullable)
\item \texttt{status} (VARCHAR, nullable)
\item \texttt{gears} (INT, nullable)
\item \texttt{license\_plate} (VARCHAR, nullable)
\item \texttt{deposit\_capacity} (DECIMAL(5,2), nullable)
\item \texttt{trunk\_capacity} (DECIMAL(5,2), nullable)
\item \texttt{horse\_power} (INT, nullable)
\item \texttt{brand\_id} (INT, foreign key to CarBrand, not null)
\item \texttt{seller\_id} (INT, foreign key to Seller, not null)
\end{itemize}

\item \textbf{Primary Key:} \texttt{id}
\item \textbf{Foreign Keys:}

\begin{itemize}
\item \texttt{brand\_id} references \texttt{CarBrand(id)}
\item \texttt{seller\_id} references \texttt{Seller(id)}
\end{itemize}

\item \textbf{Purpose:} Stores details about each car.
\end{itemize}

\subsubsection{Table: \texttt{Client}}
\label{sec:orgb9d78d6}
\begin{itemize}
\item \textbf{Attributes:}

\begin{itemize}
\item \texttt{id} (INT, primary key, auto-increment)
\item \texttt{name} (VARCHAR, not null)
\end{itemize}

\item \textbf{Primary Key:} \texttt{id}
\item \textbf{Purpose:} Stores client information.
\end{itemize}

\subsubsection{Table: \texttt{Purchase}}
\label{sec:orgbe5021f}
\begin{itemize}
\item \textbf{Attributes:}

\begin{itemize}
\item \texttt{id} (INT, primary key, auto-increment)
\item \texttt{client\_id} (INT, foreign key to Client, not null)
\item \texttt{car\_id} (INT, foreign key to Car, not null)
\item \texttt{date} (DATE, not null)
\end{itemize}

\item \textbf{Primary Key:} \texttt{id}
\item \textbf{Foreign Keys:}

\begin{itemize}
\item \texttt{client\_id} references \texttt{Client(id)}
\item \texttt{car\_id} references \texttt{Car(id)}
\end{itemize}

\item \textbf{Purpose:} Records the purchases of cars by clients.
\end{itemize}

\subsection{Notes on the Database Design}
\label{sec:org6e3dc71}
\begin{itemize}
\item \textbf{Normalization:} The database design should be normalized to at least
the third normal form (3NF) to reduce redundancy and improve data
integrity.
\item \textbf{Data Types:} Choose appropriate data types for each field. For
instance, \texttt{VARCHAR} for strings, \texttt{INT} for integers, \texttt{DECIMAL} for
prices and capacities, and \texttt{DATE} for dates.
\item \textbf{Constraints:}

\begin{itemize}
\item \texttt{NOT NULL} constraints are used for mandatory fields.
\item \texttt{AUTO\_INCREMENT} is used for primary keys to create unique
identifiers for each row.
\item \texttt{UNIQUE} constraints ensure values like brand names are not
duplicated.
\end{itemize}

\item \textbf{Foreign Key Relationships:} These enforce referential integrity
between tables. For example, each car is associated with a specific
brand and seller.
\item \textbf{Nullable Fields:} Fields like \texttt{engine\_features}, \texttt{gas\_type},
\texttt{mileage}, etc., are allowed to be \texttt{NULL} to accommodate varying
levels of detail available for different cars.
\end{itemize}

This database structure provides a comprehensive foundation for a car
dealership system, enabling efficient data storage and retrieval for a
wide range of dealership operations.

\section{PHP and HTML Code}
\label{sec:org86cff2d}
\subsection{List ALL Cars}
\label{sec:org1ccf60d}
To create a simple HTML page and its corresponding PHP script for
listing all cars with their model, brand, and price in a car dealership
project, you'll need to follow these steps. This example assumes you
have a database with the necessary tables and data in place.

\subsubsection{Step 1: Create the HTML Page}
\label{sec:org99f0a8a}
This page (\texttt{list\_cars.html}) will be simple, primarily serving as a link
to the PHP script that will display the car list.

\begin{minted}[]{html}
<!DOCTYPE html>
<html>
<head>
    <title>Car Dealership - Car List</title>
</head>
<body>
    <h1>Car List</h1>
    <p>
        <a href="list_cars.php">View All Cars</a>
    </p>
</body>
</html>
\end{minted}

\subsubsection{Step 2: Create the PHP Script}
\label{sec:org36f5a87}
This PHP script (\texttt{list\_cars.php}) will connect to the database, retrieve
the list of cars, and display the model, brand, and price. We'll use PDO
for database interaction with prepared statements.

\begin{minted}[]{php}
<?php
$host = 'localhost';
$db   = 'car_dealership';
$user = 'root';      // Replace with your database username
$pass = '';          // Replace with your database password
$charset = 'utf8mb4';

$dsn = "mysql:host=$host;dbname=$db;charset=$charset";
$options = [
    PDO::ATTR_ERRMODE            => PDO::ERRMODE_EXCEPTION,
    PDO::ATTR_DEFAULT_FETCH_MODE => PDO::FETCH_ASSOC,
    PDO::ATTR_EMULATE_PREPARES   => false,
];

try {
    $pdo = new PDO($dsn, $user, $pass, $options);
} catch (\PDOException $e) {
    throw new \PDOException($e->getMessage(), (int)$e->getCode());
}

$stmt = $pdo->prepare("SELECT Car.model, Car.price, CarBrand.name AS brand FROM Car JOIN CarBrand ON Car.brand_id = CarBrand.id");
$stmt->execute();

echo "<h1>Car List</h1>";
echo "<table border='1'><tr><th>Model</th><th>Brand</th><th>Price</th></tr>";
while ($row = $stmt->fetch()) {
    echo "<tr><td>".$row["model"]."</td><td>".$row["brand"]."</td><td>".$row["price"]."</td></tr>";
}
echo "</table>";
?>
\end{minted}

\subsubsection{Explanation:}
\label{sec:orged183c0}
\begin{enumerate}
\item \textbf{Database Connection}: The script starts by setting up a PDO
connection to the database.

\item \textbf{Prepared Statement}: A prepared statement is used to query the
database. This query joins the \texttt{Car} and \texttt{CarBrand} tables to get the
model, brand, and price of each car.

\item \textbf{Data Display}: The fetched data is displayed in an HTML table.
\end{enumerate}

\subsubsection{Note:}
\label{sec:orgdb29b83}
\begin{itemize}
\item Make sure to replace the \texttt{\$host}, \texttt{\$db}, \texttt{\$user}, and \texttt{\$pass}
variables with your actual database host, database name, username, and
password.
\item The structure of the SQL query depends on your database schema. Adjust
the field names and table names according to your actual database
schema.
\item Ensure that your PHP environment is correctly set up to connect to
your MySQL database.
\end{itemize}

This setup provides a basic and secure way to list cars from your
database using PHP and PDO with prepared statements.

\subsection{One car deatails}
\label{sec:orgaa0b000}
To create an HTML page and a corresponding PHP script for showing the
details of a specific car selected by its ID, you'll need to follow
these steps. This example assumes you have a database with the necessary
tables and data, including detailed information about each car.

\subsubsection{Step 1: Create the HTML Page for Car Selection}
\label{sec:org6b8f0cf}
This page (\texttt{select\_car.html}) allows the user to enter the ID of the car
they want to view.

\begin{minted}[]{html}
<!DOCTYPE html>
<html>
<head>
    <title>Car Dealership - Select Car</title>
</head>
<body>
    <h1>Select a Car to View Details</h1>
    <form action="car_detail.php" method="get">
        Car ID: <input type="text" name="car_id">
        <input type="submit" value="View Details">
    </form>
</body>
</html>
\end{minted}

\subsubsection{Step 2: Create the PHP Script for Displaying Car Details}
\label{sec:orgc482337}
This PHP script (\texttt{car\_detail.php}) will connect to the database,
retrieve the details of the car with the given ID, and display them.
We'll continue using PDO for database interaction with prepared
statements.

\begin{minted}[]{php}
<?php
$host = 'localhost';
$db   = 'car_dealership';
$user = 'root';      // Replace with your database username
$pass = '';          // Replace with your database password
$charset = 'utf8mb4';

$dsn = "mysql:host=$host;dbname=$db;charset=$charset";
$options = [
    PDO::ATTR_ERRMODE            => PDO::ERRMODE_EXCEPTION,
    PDO::ATTR_DEFAULT_FETCH_MODE => PDO::FETCH_ASSOC,
    PDO::ATTR_EMULATE_PREPARES   => false,
];

try {
    $pdo = new PDO($dsn, $user, $pass, $options);
} catch (\PDOException $e) {
    throw new \PDOException($e->getMessage(), (int)$e->getCode());
}

$carId = $_GET['car_id'] ?? ''; // Retrieve car ID from URL parameter

if ($carId) {
    $stmt = $pdo->prepare("SELECT * FROM Car WHERE id = :carId");
    $stmt->bindParam(':carId', $carId, PDO::PARAM_INT);
    $stmt->execute();

    $car = $stmt->fetch();
    if ($car) {
        echo "<h1>Car Details</h1>";
        echo "<p>Model: " . htmlspecialchars($car["model"]) . "</p>";
        echo "<p>Brand: " . htmlspecialchars($car["brand"]) . "</p>";
        // Display other details like price, condition, engine features, etc.
    } else {
        echo "<p>Car not found.</p>";
    }
} else {
    echo "<p>No car selected.</p>";
}
?>
\end{minted}

\subsubsection{Explanation:}
\label{sec:org6b3ba8f}
\begin{enumerate}
\item \textbf{Database Connection}: The script sets up a PDO connection to the
database.

\item \textbf{Retrieving Car ID}: The car ID is obtained from the URL parameter
(\texttt{\$\_GET['car\_id']}).

\item \textbf{Query and Display Car Details}: A prepared statement is used to
retrieve and display the details of the car with the provided ID. The
data is sanitized using \texttt{htmlspecialchars} to prevent XSS attacks.
\end{enumerate}

\subsubsection{Note:}
\label{sec:org04e5d0f}
\begin{itemize}
\item Replace the database connection details (\texttt{\$host}, \texttt{\$db}, \texttt{\$user},
\texttt{\$pass}) with your actual database information.
\item Ensure your table and column names in the SQL query match those in
your database schema.
\item The script assumes that the \texttt{Car} table contains all necessary columns
like \texttt{model}, \texttt{brand}, \texttt{price}, \texttt{condition}, etc. Modify the script to
include all relevant details as per your database schema.
\item This setup provides a basic and secure way to display detailed
information about a car from your database using PHP and PDO. For a
more user-friendly interface, you might consider adding error handling
for invalid IDs and styling the output with CSS.
\end{itemize}

\subsection{Create a new car}
\label{sec:orgd783d5c}
Creating an HTML form and corresponding PHP script for inserting a new
car with all its details into a database involves two main parts:
designing the form and handling the form submission with PHP. Below is a
basic example. This assumes you have a database structure ready to store
these details.

\subsubsection{Step 1: Create the HTML Form}
\label{sec:org7d6bf5f}
Create an HTML page (\texttt{add\_car.html}) for the input form. This form will
include various fields for car details. Some fields will be marked as
mandatory.

\begin{minted}[]{html}
<!DOCTYPE html>
<html>
<head>
    <title>Add New Car</title>
</head>
<body>
    <h1>Add New Car</h1>
    <form action="insert_car.php" method="post">
        Model (required): <input type="text" name="model" required><br>
        Brand (required): <input type="text" name="brand" required><br>
        Price (required): <input type="number" name="price" required><br>
        Condition: <input type="text" name="condition"><br>
        Engine Features: <input type="text" name="engine_features"><br>
        Gas Type: <input type="text" name="gas_type"><br>
        Mileage: <input type="number" name="mileage"><br>
        Status (new/refurbished): <input type="text" name="status"><br>
        Gears: <input type="number" name="gears"><br>
        License Plate: <input type="text" name="license_plate"><br>
        Deposit Capacity: <input type="number" name="deposit_capacity"><br>
        Trunk Capacity: <input type="number" name="trunk_capacity"><br>
        Horse Power: <input type="number" name="horse_power"><br>
        <input type="submit" value="Add Car">
    </form>
</body>
</html>
\end{minted}

\subsubsection{Step 2: Create the PHP Script for Form Handling}
\label{sec:org17903ce}
Create a PHP script (\texttt{insert\_car.php}) to handle the form submission.
This script will insert the new car details into your database.

\begin{minted}[]{php}
<?php
$host = 'localhost';
$db   = 'car_dealership';
$user = 'root';      // Replace with your database username
$pass = '';          // Replace with your database password
$charset = 'utf8mb4';

$dsn = "mysql:host=$host;dbname=$db;charset=$charset";
$options = [
    PDO::ATTR_ERRMODE            => PDO::ERRMODE_EXCEPTION,
    PDO::ATTR_DEFAULT_FETCH_MODE => PDO::FETCH_ASSOC,
    PDO::ATTR_EMULATE_PREPARES   => false,
];

try {
    $pdo = new PDO($dsn, $user, $pass, $options);
} catch (\PDOException $e) {
    throw new \PDOException($e->getMessage(), (int)$e->getCode());
}

if ($_SERVER["REQUEST_METHOD"] == "POST") {
    $stmt = $pdo->prepare("INSERT INTO Car (model, brand, price, condition, engine_features, gas_type, mileage, status, gears, license_plate, deposit_capacity, trunk_capacity, horse_power) VALUES (?, ?, ?, ?, ?, ?, ?, ?, ?, ?, ?, ?, ?)");

    $stmt->execute([
        $_POST['model'],
        $_POST['brand'],
        $_POST['price'],
        $_POST['condition'] ?? null,
        $_POST['engine_features'] ?? null,
        $_POST['gas_type'] ?? null,
        $_POST['mileage'] ?? null,
        $_POST['status'] ?? null,
        $_POST['gears'] ?? null,
        $_POST['license_plate'] ?? null,
        $_POST['deposit_capacity'] ?? null,
        $_POST['trunk_capacity'] ?? null,
        $_POST['horse_power'] ?? null
    ]);

    echo "New car added successfully.";
}
?>
\end{minted}

\subsubsection{Explanation}
\label{sec:org02c408b}
\begin{itemize}
\item The HTML form includes various fields for entering car details. The
\texttt{required} attribute is added to mandatory fields.
\item The PHP script connects to the database and uses a prepared statement
to insert the car data. It uses the POST method to retrieve data from
the form.
\item Null coalescing operators (\texttt{??}) are used for optional fields to
handle cases where no data is entered.
\item Proper error handling should be implemented for a production
environment (e.g., try-catch blocks, validation of input data).
\end{itemize}

\subsubsection{Note}
\label{sec:org959a493}
\begin{itemize}
\item Replace the \texttt{\$host}, \texttt{\$db}, \texttt{\$user}, and \texttt{\$pass} variables with your
actual database host, database name, username, and password.
\item The SQL query in the prepared statement should match your database's
table structure. Adjust the field names and table name according to
your database schema.
\item This example does not include client-side or server-side validation
for the sake of brevity. In a real-world application, you should
validate and sanitize all inputs to ensure data integrity and
security.
\item The form assumes that the user knows the Car ID they wish to update.
In a practical scenario, you might have a selection process or a
search feature to find the car before updating its details.
\end{itemize}
\end{document}
